\documentclass[11pt,a4paper]{article}

% --- Configuración de Idioma y Tipografía Nativa ---
\usepackage[utf8]{inputenc}
\usepackage[spanish,es-tabla]{babel}
\usepackage[T1]{fontenc}

% --- Geometría y Configuración de Página ---
\usepackage[top=2.5cm, bottom=2.5cm, left=2.5cm, right=2.5cm]{geometry}
\usepackage{parskip} % Espaciado automático limpio entre párrafos sin sangría

% --- Paquetes para Gráficos, Matemáticas y Tablas Profesionales ---
\usepackage{graphicx}
\usepackage{booktabs}
\usepackage{tabularx}
\usepackage{array}
\usepackage{amsmath}
\usepackage{amsfonts}
\usepackage{hyperref}

% Definición de columnas automáticas con alineación a la izquierda
\newcolumntype{L}{>{\raggedright\arraybackslash}X}

\hypersetup{
    colorlinks=true,
    linkcolor=blue,
    filecolor=magenta,      
    urlcolor=cyan,
}

\begin{document}

% ==============================================================================
% PORTADA
% ==============================================================================
\begin{titlepage}
    \centering
    \vspace*{2cm}
    
    {\Huge\bfseries INFORME TÉCNICO}\\[0.6cm]
    {\large\bfseries AP13 · Propagación del impacto por dependencias}\\[0.4cm]
    {\Large\bfseries Fortune Bank}\\[2cm]
    
    \begin{center}
        \textbf{Alumno:} Xavier Margalef Riestra\\[0.3cm]
        \textbf{Metodología:} MAGERIT v3 · Esquema Nacional de Seguridad (ENS)\\[0.3cm]
        \textbf{Actividad:} Dependencias entre activos y propagación del impacto\\[0.3cm]
        \textbf{Fecha:} Julio 2026
    \end{center}
    
    \vfill
    
    \includegraphics[width=0.5\textwidth]{image-2.png}
    
    \vspace*{2cm}
\end{titlepage}

\pagebreak

% ==============================================================================
% INTRODUCCIÓN Y PLANTEAMIENTO
% ==============================================================================
\section*{Introducción}
Este informe analiza la propagación del impacto entre activos bajo el modelo de la metodología \textbf{MAGERIT v3}. El objetivo es determinar cómo una degradación operativa en un activo de soporte escala verticalmente afectando a los componentes superiores, calculando el impacto final transferido al servicio de negocio. 

Asimismo, se evalúa la eficacia de la implementación de \textbf{salvaguardas} para la obtención del impacto residual y la determinación del patrón estructural de dependencias de la arquitectura.

\section*{Planteamiento}
Evaluación del comportamiento del sistema de transferencias de \textbf{Fortune Bank} ante una caída por fallo de hardware.

La infraestructura consta de los siguientes componentes:
\begin{itemize}
    \item \textbf{Activo A:} Sistema Core de transferencias.
    \item \textbf{Nodo B:} Middleware de procesamiento.
    \item \textbf{Nodo C:} Middleware de procesamiento.
    \item \textbf{Servidor D:} Servidor hardware base del que dependen ambos middlewares.
\end{itemize}

\textbf{Métricas del escenario:}
\begin{itemize}
    \item El \textbf{Servidor D} sufre una degradación inicial del \textbf{50 \%} por fallo físico.
    \item Los nodos \textbf{B} y \textbf{C} tienen un grado de dependencia del \textbf{70 \%} respecto al Servidor D.
    \item El \textbf{Sistema Core (A)} depende en un \textbf{60 \%} de cada uno de los middlewares de forma independiente.
    \item Las salvaguardas aplicadas presentan una eficacia frente al impacto del 50 \%.
\end{itemize}

\pagebreak

% ==============================================================================
% IDENTIFICACIÓN DE ACTIVOS Y TOPOLOGÍA
% ==============================================================================
\section*{Identificación de activos}
Aplicando los criterios de \textbf{MAGERIT v3}, se clasifican los activos diferenciando los elementos esenciales de negocio de los activos de soporte de infraestructura.

\begin{table}[h!]
    \small
    \centering
    \begin{tabularx}{\textwidth}{l l L}
        \toprule
        \textbf{Activo} & \textbf{Tipo} & \textbf{Función} \\
        \midrule
        \textbf{Activo A} & Activo esencial & Sistema Core de transferencias. Servicio crítico de negocio. \\
        \textbf{Nodo B} & Activo de soporte & \href{https://owasp.org/www-community/Application_Architecture}{Middleware} de procesamiento transaccional. \\
        \textbf{Nodo C} & Activo de soporte & Middleware homólogo al Nodo B. \\
        \textbf{Servidor D} & Soporte crítico & Servidor hardware compartido. Constituye un \href{https://en.wikipedia.org/wiki/Single_point_of_failure}{Single Point of Failure (SPOF)}. \\
        \bottomrule
    \end{tabularx}
\end{table}

Dado que el valor del negocio reside en el activo esencial, este se transmite de forma descendente. Consiguientemente, un incidente en el componente base \textbf{Servidor D} repercute de forma ascendente directa en la disponibilidad operativa del \textbf{Sistema Core}.

\begin{figure}[h!]
    \centering
    \includegraphics[width=0.6\textwidth]{image.png}
\end{figure}

\section*{Topología de dependencias}
El flujo jerárquico de dependencias se estructura de la siguiente manera:

\begin{verbatim}
          Activo A
         /        \
      60 %      60 %
      /            \
 Nodo B          Nodo C
      \            /
      70 %      70 %
         \      /
        Servidor D
\end{verbatim}

La arquitectura define una \textbf{dependencia convergente} (topología en diamante). El \textbf{Servidor D} actúa como un \textbf{SPOF}: un impacto en este componente compromete simultáneamente ambos canales intermedios, al carecer el diseño actual de mecanismos de tolerancia a fallos o balanceo de carga.

\pagebreak

% ==============================================================================
% DESARROLLO MATEMÁTICO - IMPACTO INICIAL Y RESIDUAL
% ==============================================================================
\section*{1. Impacto inicial sobre el Sistema Core (Activo A)}
Ante una degradación del 50 \% en el Servidor D, el cálculo del impacto repercutido se basa en la fórmula formal de la metodología:
\[ \text{Impacto Repercutido} = \text{Degradación} \cdot \text{Grado de Dependencia} \]

\subsection*{Paso 1: Del Servidor D a los Nodos intermedios de soporte}
\begin{itemize}
    \item \textbf{Nodo B:} \[ 0,50 \cdot 0,70 = 0,35 \rightarrow \mathbf{35\ \%\ \text{de impacto}} \]
    \item \textbf{Nodo C:} \[ 0,50 \cdot 0,70 = 0,35 \rightarrow \mathbf{35\ \%\ \text{de impacto}} \]
\end{itemize}

\subsection*{Paso 2: De los Nodos intermedios al Activo Esencial (A)}
Cada canal de procesamiento traslada el impacto acumulado hacia el Activo A aplicando su grado de dependencia correspondiente del \textbf{60 \%} ($0,60$):
\begin{itemize}
    \item \textbf{Ruta B $\rightarrow$ A:} \[ 0,35 \cdot 0,60 = 0,21 \rightarrow \mathbf{21\ \%\ \text{de impacto}} \]
    \item \textbf{Ruta C $\rightarrow$ A:} \[ 0,35 \cdot 0,60 = 0,21 \rightarrow \mathbf{21\ \%\ \text{de impacto}} \]
\end{itemize}

\subsection*{Cálculo del Impacto Total Acumulado}
Al coexistir dos rutas de transmisión convergentes simultáneas sin redundancia activa en esta fase, se agrega el impacto recibido por ambas vías sobre el activo de negocio:
\[ 21\% + 21\% = 42\% \]
\textbf{Impacto inicial total en el Sistema Core (Activo A): 42 \%.}

\section*{2. Impacto residual tras implantar salvaguardas}
Se introducen medidas orientadas a acotar la degradación (implementación de cachés locales en los Nodos B y C) con una eficacia frente al impacto ($e_i$) del \textbf{50 \%} ($0,50$). Aplicando la ecuación matemática de impacto residual:
\[ \text{Impacto Residual} = \text{Impacto Inicial} \cdot (1 - e_i) \]

\begin{itemize}
    \item \textbf{Ruta B mitigada:} \[ 21\% \cdot (1 - 0,50) = 10,5\% \]
    \item \textbf{Ruta C mitigada:} \[ 21\% \cdot (1 - 0,50) = 10,5\% \]
\end{itemize}
\textbf{Impacto residual final consolidado en el Core:} \[ 10,5\% + 10,5\% = \mathbf{21\%} \]

\pagebreak

% ==============================================================================
% MATRIZ DE RESUMEN Y PATRÓN
% ==============================================================================
\section*{3. Patrón de dependencias y resumen}

\begin{figure}[h!]
    \centering
    \includegraphics[width=0.75\textwidth]{image-1.png}
\end{figure}

La expresión formal que define la topología estructural en diamante corresponde a:
\[ (A \rightarrow B) \land (A \rightarrow C) \land (B \rightarrow D) \land (C \rightarrow D) \]

\begin{table}[h!]
    \small
    \centering
    \begin{tabularx}{\textwidth}{L r}
        \toprule
        \textbf{Concepto Técnico} & \textbf{Valor Operativo} \\
        \midrule
        Degradación inicial del Servidor D ($d$) & \textbf{50 \%} \\
        Grado de Dependencia Nodos (B/C) $\rightarrow$ Servidor D & \textbf{70 \%} ($0,70$) \\
        Impacto intermedio en Nodos B y C & \textbf{35 \%} \\
        Grado de Dependencia Activo A $\rightarrow$ Nodos (B/C) & \textbf{60 \%} ($0,60$) \\
        \midrule
        \textbf{Impacto inicial total en Activo A} & \textbf{42 \%} \\
        \textbf{Impacto residual total en Activo A} & \textbf{21 \%} \\
        \midrule
        Patrón estructural de la arquitectura & \textbf{Dependencia en Diamante} \\
        Componente crítico detectado & \textbf{Servidor D (SPOF)} \\
        \bottomrule
    \end{tabularx}
\end{table}

\pagebreak

% ==============================================================================
% ESCENARIOS ALTERNATIVOS (INGENIERÍA INVERSA)
% ==============================================================================
\section*{Escenarios alternativos (Ingeniería Inversa)}

Para evaluar los escenarios alternativos se realiza un análisis de ingeniería inversa. Partiendo del impacto final reflejado en el activo raíz ($i_A$), se realiza el despeje secuencial descendente para identificar la degradación requerida en el activo base de infraestructura.

Las expresiones matemáticas utilizadas para el despeje directo son:
\begin{enumerate}
    \item \textbf{Impacto intermedio en el nivel de Middleware:} 
    \[ \text{Impacto}_{\text{Nodo}} = \frac{\text{Impacto}_{\text{Ruta}}}{\text{grado}(A \rightarrow \text{Nodo})} \]
    \item \textbf{Degradación requerida en el componente hardware raíz (Servidor D):} 
    \[ \text{Degradación}_D = \frac{\text{Impacto}_{\text{Nodo}}}{\text{grado}(\text{Nodo} \rightarrow D)} \]
\end{enumerate}
\textit{Constantes fijas del sistema:} $\text{grado}(A \rightarrow B/C) = 0,60$ | $\text{grado}(B/C \rightarrow D) = 0,70$ | Eficacia de salvaguarda ($e_i$) = $50\%$.

\vspace{0.3cm}

\begin{table}[h!]
    \footnotesize
    \centering
    \begin{tabularx}{\textwidth}{p{2.3cm} c L c c c}
        \toprule
        \textbf{Escenario} & \textbf{\bfseries $i_A$} & \textbf{Cálculo de Ingeniería Inversa (*Top-Down*)} & \textbf{Nodos} & \textbf{\bfseries $d_D$} & \textbf{\bfseries Residuo} \\
        \midrule
        \textbf{1. Diamante} & 42\% & 1. Por rama: $\text{Impacto}_{\text{Ruta}} = 21\%$.\newline 2. En nodo: $\text{Impacto}_B = \frac{0,21}{0,60} = 0,35$.\newline 3. En raíz: $\text{Degradación}_D = \frac{0,35}{0,70} = 0,50$. & \begin{tabular}{@{}c@{}}$d_B=35\%$\\$d_C=35\%$\end{tabular} & \textbf{50\%} & \textbf{21\%} \\
        \midrule
        \textbf{2. Ruta Única} & 21\% & 1. Un solo canal activo transmite: $\text{Impacto}_{\text{Ruta}} = 21\%$.\newline 2. En nodo: $\text{Impacto}_B = \frac{0,21}{0,60} = 0,35$.\newline 3. En raíz: $\text{Degradación}_D = \frac{0,35}{0,70} = 0,50$. & \begin{tabular}{@{}c@{}}$d_B=35\%$\\$d_C=0\%$\end{tabular} & \textbf{50\%} & \textbf{10,5\%} \\
        \midrule
        \textbf{3. Conexión Directa} & 35\% & 1. Arquitectura plana. A acopla directo a D ($\text{grado} = 0,70$).\newline 2. En raíz: $\text{Degradación}_D = \frac{0,35}{0,70} = 0,50$. & \begin{tabular}{@{}c@{}}$d_B=\text{N/A}$\\$d_C=\text{N/A}$\end{tabular} & \textbf{50\%} & \textbf{17,5\%} \\
        \midrule
        \textbf{4. Doble Salvaguarda} & 21\% & 1. Flujo idéntico de canal al Escenario 2.\newline 2. $\text{Impacto}_B = 35\% \rightarrow \text{Degradación}_D = 50\%$.\newline 3. Controles en serie: $21\% \cdot (1-0,50) \cdot (1-0,50)$. & \begin{tabular}{@{}c@{}}$d_B=35\%$\\$d_C=0\%$\end{tabular} & \textbf{50\%} & \textbf{5,25\%} \\
        \bottomrule
    \end{tabularx}
\end{table}

\begin{figure}[h!]
    \centering
    \includegraphics[width=0.75\textwidth]{image-3.png}
\end{figure}

\pagebreak

% ==============================================================================
% REFERENCIAS
% ==============================================================================
\section*{Referencias}
\subsection*{Metodologías y Normativa Oficial}
\begin{itemize}
    \item \textbf{MAGERIT v3} -- Metodología de Análisis y Gestión de Riesgos de los Sistemas de Información. Ministerio de Asuntos Económicos y Transformación Digital.
    \item \textbf{Esquema Nacional de Seguridad (ENS)} -- Real Decreto 311/2022. Centro Criptológico Nacional (CCN).
    \item \textbf{ISO/IEC 27001:2022} -- Information Security Management Systems -- Requirements.
\end{itemize}

\subsection*{Marcos Técnicos}
\begin{itemize}
    \item \textbf{Single Point of Failure (SPOF)} -- OWASP Application Architecture Mitigation Standards.
    \item \textbf{Defense in Depth} -- NIST Cybersecurity Framework (CSF 2.0).
\end{itemize}

\end{document}